\documentclass{article}
\usepackage[utf8]{inputenc}
\usepackage[T1]{fontenc}
\usepackage{geometry}
\usepackage{graphicx}
\usepackage{float}
\usepackage{hyperref}
\usepackage[english,indonesian]{babel}

\title{Optimasi Pemrosesan Template LaTeX Menggunakan Pendekatan Berbasis AI}
\author{Academic Support Agent}
\date{\today}

\begin{document}
\maketitle

\begin{abstract}

\end{abstract}

\chapter{PENDAHULUAN}


\chapter{TINJAUAN PUSTAKA}


% Original content (commented out):
% \documentclass[12pt]{article}
% Þetta lokaverkefnis template er skrifað af Jóhannes Bergi og uppfært þann 18. ágúst 2021 síðast. 
% Hugsað fyrst of fremst fyrir lokaverkefni í tæknifræði en líklega almennt fyrir skýrslur og fleira gæti þetta verið ágætt. 
% Ég henti lykilsida.tex nokkuð gróflega saman þegar ég gerði lokaverkefnið, ég hef ekki haft tíma til þess að gera kóðan læsilegri og fallegri, en það er klárlega á stefnuskránni.
%Logo á forsíðu er gert í TikZ og litir fengnir frá síðu HR.
\usepackage[margin=2.54cm,includefoot, headheight=110pt]{geometry}

\input{Main_pakkar}
\input{Matlab/matlab_forpyttur}


\usepackage[colorlinks=true, linkcolor=black, citecolor=blue, urlcolor=blue, bookmarks=true, breaklinks=true]{hyperref} % always last package! 
\usepackage{bookmark} % Bookmark is extention of hyperref, so I guess this goes last.

\begin{document}
\pagenumbering{roman}
\input{kaflar/Forsida.tex}\setcounter{page}{1}
\input{kaflar/lykilsida.tex}
\newpage

\input{kaflar/formali}
\newpage
\tableofcontents
\newpage
\listoffigures \addcontentsline{toc}{section}{List of Figures}
\listoftables \addcontentsline{toc}{section}{List of Tables}
\newpage
\pagenumbering{arabic}
\input{kaflar/Introduction}
\input{kaflar/mainmatter}

\newpage
\bibliographystyle{ieeetr}
\bibliography{bibliography.bib}\addcontentsline{toc}{section}{References}

\newpage
\addcontentsline{toc}{section}{Appendices}
\input{kaflar/Vidauki}
		
\end{document}


\end{document}